% Fichier engendré par outils/generer-tex.py à partir de EntiersDivisibilite.lean.
% Ne pas éditer à la main : tout le texte provient des commentaires du fichier Lean.
\documentclass[11pt,a4paper]{article}

% Compile aussi bien avec pdflatex qu'avec xelatex ou tectonic.
\usepackage{iftex}
\ifPDFTeX
  \usepackage[T1]{fontenc}
  \usepackage[utf8]{inputenc}
\else
  \usepackage{fontspec}
\fi
\usepackage[french]{babel}
\usepackage{amsmath,amssymb,amsthm}
\usepackage[margin=2.5cm]{geometry}
\usepackage[hidelinks]{hyperref}

\theoremstyle{definition}
\newtheorem{definition}{Définition}
\theoremstyle{plain}
\newtheorem{theoreme}{Théorème}
\newtheorem{lemme}[theoreme]{Lemme}

% Renvoi à la déclaration Lean, une ligne après la démonstration, aligné à droite.
\newcommand{\source}[2]{\par\smallskip\noindent\hfill{\footnotesize\href{#1}{\texttt{#2}}}\par\smallskip}

\title{\texttt{EntiersDivisibilite}}
\date{}

\begin{document}
\maketitle
% Les identifiants Lean en machine à écrire ne se coupent pas : on tolère des
% espacements plus lâches plutôt que des lignes qui débordent.
\sloppy

\noindent
Collège \textemdash{} parité, divisibilité, division euclidienne, nombres premiers et PGCD.
Reprend les quatorze énoncés de la section \og{} Entiers, divisibilité \fg{} du chapitre.

Lean core seulement : pas d'import de Mathlib. \texttt{Pair}, \texttt{Impair}, \texttt{Premier} et la somme
des chiffres sont donc définis ici plutôt que repris d'une bibliothèque.

\begin{definition}
\texttt{n} est \emph{pair} s'il existe un entier naturel $k$ tel que $n = 2k$.
\end{definition}

\source{https://github.com/Commutator-IO/learn-lean/blob/main/cours/01-college/07-enonces-conserves/EntiersDivisibilite.lean\#L9}{EntiersDivisibilite.lean\#L9}

\begin{definition}
\texttt{n} est \emph{impair} s'il existe un entier naturel $k$ tel que $n = 2k + 1$.
\end{definition}

\source{https://github.com/Commutator-IO/learn-lean/blob/main/cours/01-college/07-enonces-conserves/EntiersDivisibilite.lean\#L12}{EntiersDivisibilite.lean\#L12}

\begin{definition}
\texttt{p} est \emph{premier} si $2 \le p$ et si tout diviseur $d$ de $p$ vérifie $d = 1$
ou $d = p$.
\end{definition}

\source{https://github.com/Commutator-IO/learn-lean/blob/main/cours/01-college/07-enonces-conserves/EntiersDivisibilite.lean\#L15}{EntiersDivisibilite.lean\#L15}

\section{Un entier est pair ou impair, jamais les deux}

\begin{theoreme}
Un entier est pair ou impair.
\end{theoreme}

\begin{proof}
Le reste $n \bmod 2$ ne peut valoir que $0$ ou $1$ ; distinguons ces deux cas. Dans
l'un comme dans l'autre, il suffit de prendre $k = \lfloor n/2 \rfloor$ : de
$n = 2\lfloor n/2 \rfloor + (n \bmod 2)$ on tire $n = 2k$ si le reste est nul, et
$n = 2k + 1$ sinon. Les deux égalités se vérifient par un calcul.
\end{proof}

\source{https://github.com/Commutator-IO/learn-lean/blob/main/cours/01-college/07-enonces-conserves/EntiersDivisibilite.lean\#L20}{EntiersDivisibilite.lean\#L20}

\begin{theoreme}
Un entier n'est pas les deux à la fois.
\end{theoreme}

\begin{proof}
Par l'absurde : si $n$ était à la fois pair et impair, on écrirait $n = 2k$ et
$n = 2l + 1$, d'où $2k = 2l + 1$. Or un nombre pair ne peut pas être égal à un nombre
impair.
\end{proof}

\source{https://github.com/Commutator-IO/learn-lean/blob/main/cours/01-college/07-enonces-conserves/EntiersDivisibilite.lean\#L26}{EntiersDivisibilite.lean\#L26}

\begin{theoreme}
Un entier est pair ou impair, jamais les deux.
\end{theoreme}

\begin{proof}
C'est la conjonction des deux théorèmes précédents.
\end{proof}

\source{https://github.com/Commutator-IO/learn-lean/blob/main/cours/01-college/07-enonces-conserves/EntiersDivisibilite.lean\#L31}{EntiersDivisibilite.lean\#L31}

\section{Somme de deux pairs = pair ; pair + impair = impair ; impair + impair = pair}

\begin{theoreme}
Somme de deux pairs = pair.
\end{theoreme}

\begin{proof}
On écrit les hypothèses sous la forme $m = 2k$ et $n = 2l$. Alors
$m + n = 2k + 2l = 2(k+l)$, et $k + l$ convient.
\end{proof}

\source{https://github.com/Commutator-IO/learn-lean/blob/main/cours/01-college/07-enonces-conserves/EntiersDivisibilite.lean\#L38}{EntiersDivisibilite.lean\#L38}

\begin{theoreme}
Pair + impair = impair.
\end{theoreme}

\begin{proof}
On écrit $m = 2k$ et $n = 2l + 1$. Alors $m + n = 2(k+l) + 1$ : on prend encore
$k + l$, et le $+1$ de $n$ devient celui de la conclusion.
\end{proof}

\source{https://github.com/Commutator-IO/learn-lean/blob/main/cours/01-college/07-enonces-conserves/EntiersDivisibilite.lean\#L45}{EntiersDivisibilite.lean\#L45}

\begin{theoreme}
Impair + impair = pair.
\end{theoreme}

\begin{proof}
On écrit $m = 2k + 1$ et $n = 2l + 1$. Alors $m + n = 2k + 2l + 2 = 2(k + l + 1)$ : les
deux $+1$ se recombinent en un facteur $2$, et $k + l + 1$ convient.
\end{proof}

\source{https://github.com/Commutator-IO/learn-lean/blob/main/cours/01-college/07-enonces-conserves/EntiersDivisibilite.lean\#L52}{EntiersDivisibilite.lean\#L52}

\section{Notion de multiple et de diviseur ; un diviseur de \texttt{n} est inférieur ou égal à \texttt{n}}

\begin{theoreme}
Notion de multiple et de diviseur.
\end{theoreme}

\begin{proof}
Les deux membres disent la même chose : sur $\mathbb{N}$, « $a$ divise $n$ » signifie
précisément « il existe un entier $k$ tel que $n = ak$ ». Il n'y a rien à vérifier.
\end{proof}

\source{https://github.com/Commutator-IO/learn-lean/blob/main/cours/01-college/07-enonces-conserves/EntiersDivisibilite.lean\#L61}{EntiersDivisibilite.lean\#L61}

\begin{theoreme}
Un diviseur de \texttt{n} est inférieur ou égal à \texttt{n}.
\end{theoreme}

\begin{proof}
C'est \texttt{Nat.le\_of\_dvd} de la bibliothèque standard : si $n = ak$ avec $n > 0$,
alors $k \ge 1$, donc $a \le ak = n$.
\end{proof}

\source{https://github.com/Commutator-IO/learn-lean/blob/main/cours/01-college/07-enonces-conserves/EntiersDivisibilite.lean\#L64}{EntiersDivisibilite.lean\#L64}

\begin{theoreme}
Contre-exemple pour \texttt{n = 0}.
\end{theoreme}

\begin{proof}
Il suffit de prendre $k = 0$, puisque $0 = a \cdot 0$.

Cet énoncé est la raison d'être de l'hypothèse $0 < n$ du théorème précédent : les
diviseurs de $0$ ne sont pas bornés. L'énoncé de la classe, « un diviseur de $n$ est
inférieur ou égal à $n$ », est donc faux tel quel ; il n'est vrai que parce qu'au
collège $n$ compte toujours des objets à partager.
\end{proof}

\source{https://github.com/Commutator-IO/learn-lean/blob/main/cours/01-college/07-enonces-conserves/EntiersDivisibilite.lean\#L68}{EntiersDivisibilite.lean\#L68}

\section{Critère de divisibilité par 2, par 5, par 10 (chiffre des unités)}

\begin{definition}
Le \emph{chiffre des unités} de $n$ est $n \bmod 10$.
\end{definition}

\source{https://github.com/Commutator-IO/learn-lean/blob/main/cours/01-college/07-enonces-conserves/EntiersDivisibilite.lean\#L74}{EntiersDivisibilite.lean\#L74}

\begin{theoreme}
Critère de divisibilité par 2 (chiffre des unités).
\end{theoreme}

\begin{proof}
Après dépliage de la définition du chiffre des unités, il s'agit de
$2 \mid n \iff 2 \mid (n \bmod 10)$, prouvée dans les deux sens : comme $2 \mid 10$, écrire $n = 10q + (n \bmod 10)$ montre que $n$ et
$n \bmod 10$ ont le même reste modulo $2$.
\end{proof}

\source{https://github.com/Commutator-IO/learn-lean/blob/main/cours/01-college/07-enonces-conserves/EntiersDivisibilite.lean\#L77}{EntiersDivisibilite.lean\#L77}

\begin{theoreme}
Critère de divisibilité par 2, chiffres énumérés.
\end{theoreme}

\begin{proof}
Même argument, la condition $2 \mid (n \bmod 10)$ étant énumérée : le reste modulo $10$
est l'un des dix chiffres, et les pairs sont $0$, $2$, $4$, $6$, $8$.
\end{proof}

\source{https://github.com/Commutator-IO/learn-lean/blob/main/cours/01-college/07-enonces-conserves/EntiersDivisibilite.lean\#L82}{EntiersDivisibilite.lean\#L82}

\begin{theoreme}
Critère de divisibilité par 5 (chiffre des unités).
\end{theoreme}

\begin{proof}
Identique, avec $5 \mid 10$ : $n$ et $n \bmod 10$ ont même reste modulo $5$, et les
chiffres divisibles par $5$ sont $0$ et $5$.
\end{proof}

\source{https://github.com/Commutator-IO/learn-lean/blob/main/cours/01-college/07-enonces-conserves/EntiersDivisibilite.lean\#L89}{EntiersDivisibilite.lean\#L89}

\begin{theoreme}
Critère de divisibilité par 10 (chiffre des unités).
\end{theoreme}

\begin{proof}
Cas limite du même argument : $n \bmod 10$ est le reste de $n$ modulo $10$, donc il est
nul exactement quand $10$ divise $n$.
\end{proof}

\source{https://github.com/Commutator-IO/learn-lean/blob/main/cours/01-college/07-enonces-conserves/EntiersDivisibilite.lean\#L95}{EntiersDivisibilite.lean\#L95}

\section{Critère de divisibilité par 3 et par 9 (somme des chiffres)}

\begin{definition}
La \emph{somme des chiffres} $S(n)$ est définie par récursion : $S(0) = 0$ et, pour
$n \ne 0$, $S(n) = (n \bmod 10) + S(\lfloor n/10 \rfloor)$. La récursion termine car
$\lfloor n/10 \rfloor < n$ dès que $n \ne 0$.
\end{definition}

\source{https://github.com/Commutator-IO/learn-lean/blob/main/cours/01-college/07-enonces-conserves/EntiersDivisibilite.lean\#L102}{EntiersDivisibilite.lean\#L102}

\begin{theoreme}
\texttt{n} et la somme de ses chiffres ont le même reste modulo 9.
\end{theoreme}

\begin{proof}
Par récurrence forte sur $n$. Si $n = 0$, les deux membres valent $0$.

Sinon, on déplie $S(n) = (n \bmod 10) + S(\lfloor n/10 \rfloor)$. Le quotient
$\lfloor n/10 \rfloor$ est strictement plus petit que $n$, ce qui autorise l'hypothèse de
récurrence : $\lfloor n/10 \rfloor \equiv S(\lfloor n/10 \rfloor) \pmod 9$. Comme
$n = 10\lfloor n/10 \rfloor + (n \bmod 10)$ et $10 \equiv 1 \pmod 9$, on obtient
\[
n \equiv \lfloor n/10 \rfloor + (n \bmod 10)
  \equiv S(\lfloor n/10 \rfloor) + (n \bmod 10) = S(n) \pmod 9 .
\]
\end{proof}

\source{https://github.com/Commutator-IO/learn-lean/blob/main/cours/01-college/07-enonces-conserves/EntiersDivisibilite.lean\#L108}{EntiersDivisibilite.lean\#L108}

\begin{theoreme}
\texttt{n} et la somme de ses chiffres ont le même reste modulo 3.
\end{theoreme}

\begin{proof}
Même récurrence, en remplaçant $9$ par $3$ : $10 \equiv 1 \pmod 3$ également.
\end{proof}

\source{https://github.com/Commutator-IO/learn-lean/blob/main/cours/01-college/07-enonces-conserves/EntiersDivisibilite.lean\#L120}{EntiersDivisibilite.lean\#L120}

\begin{theoreme}
Critère de divisibilité par 3 (somme des chiffres).
\end{theoreme}

\begin{proof}
Conséquence immédiate du lemme : $n$ et $S(n)$ ont le même reste modulo $3$, donc l'un
est divisible par $3$ si et seulement si l'autre l'est.
\end{proof}

\source{https://github.com/Commutator-IO/learn-lean/blob/main/cours/01-college/07-enonces-conserves/EntiersDivisibilite.lean\#L132}{EntiersDivisibilite.lean\#L132}

\begin{theoreme}
Critère de divisibilité par 9 (somme des chiffres).
\end{theoreme}

\begin{proof}
Identique avec le lemme modulo $9$.
\end{proof}

\source{https://github.com/Commutator-IO/learn-lean/blob/main/cours/01-college/07-enonces-conserves/EntiersDivisibilite.lean\#L137}{EntiersDivisibilite.lean\#L137}

\section{Critère de divisibilité par 4 (deux derniers chiffres)}

\begin{definition}
Le nombre formé par les \emph{deux derniers chiffres} de $n$ est $n \bmod 100$.
\end{definition}

\source{https://github.com/Commutator-IO/learn-lean/blob/main/cours/01-college/07-enonces-conserves/EntiersDivisibilite.lean\#L144}{EntiersDivisibilite.lean\#L144}

\begin{theoreme}
Critère de divisibilité par 4 (deux derniers chiffres).
\end{theoreme}

\begin{proof}
Même schéma que pour $2$ et $5$, cette fois avec $4 \mid 100$ : écrire
$n = 100q + (n \bmod 100)$ montre que $n$ et $n \bmod 100$ ont même reste modulo $4$.
\end{proof}

\source{https://github.com/Commutator-IO/learn-lean/blob/main/cours/01-college/07-enonces-conserves/EntiersDivisibilite.lean\#L147}{EntiersDivisibilite.lean\#L147}

\section{Si \texttt{a \(\mid\) b} et \texttt{a \(\mid\) c} alors \texttt{a \(\mid\) (b + c)} et \texttt{a \(\mid\) (b − c)}}

\begin{theoreme}
Si \texttt{a \(\mid\) b} et \texttt{a \(\mid\) c} alors \texttt{a \(\mid\) (b + c)}.
\end{theoreme}

\begin{proof}
Si $b = ak$ et $c = al$, alors $b + c = a(k+l)$ (\texttt{Nat.dvd\_add}).
\end{proof}

\source{https://github.com/Commutator-IO/learn-lean/blob/main/cours/01-college/07-enonces-conserves/EntiersDivisibilite.lean\#L154}{EntiersDivisibilite.lean\#L154}

\begin{theoreme}
Si \texttt{a \(\mid\) b} et \texttt{a \(\mid\) c} alors \texttt{a \(\mid\) (b − c)}.
\end{theoreme}

\begin{proof}
\texttt{Nat.dvd\_sub}. Sur $\mathbb{N}$, la soustraction est tronquée : si $b < c$, la
différence vaut $0$, que tout entier divise, et l'énoncé reste vrai. L'énoncé scolaire,
lui, suppose implicitement $c \le b$, sans quoi la différence n'aurait pas de sens.
\end{proof}

\source{https://github.com/Commutator-IO/learn-lean/blob/main/cours/01-college/07-enonces-conserves/EntiersDivisibilite.lean\#L158}{EntiersDivisibilite.lean\#L158}

\section{Division euclidienne : existence et unicité de \texttt{(q, r)} avec \texttt{a = bq + r}, \texttt{0 \(\le\) r < b}}

\begin{theoreme}
Division euclidienne : existence de \texttt{(q, r)} avec \texttt{a = bq + r}, \texttt{0 \(\le\) r < b}.
\end{theoreme}

\begin{proof}
Prendre $q = \lfloor a/b \rfloor$ et $r = a \bmod b$. L'égalité $a = bq + r$ est
\texttt{Nat.div\_add\_mod}, et l'inégalité $r < b$ est \texttt{Nat.mod\_lt}, qui demande
$b > 0$.
\end{proof}

\source{https://github.com/Commutator-IO/learn-lean/blob/main/cours/01-college/07-enonces-conserves/EntiersDivisibilite.lean\#L164}{EntiersDivisibilite.lean\#L164}

\begin{theoreme}
Division euclidienne : unicité de \texttt{(q, r)}.
\end{theoreme}

\begin{proof}
Le critère \texttt{Nat.div\_mod\_unique} dit que, pour $b > 0$, la conjonction de
$r + bq = a$ et $r < b$ caractérise $q = \lfloor a/b \rfloor$ et $r = a \bmod b$.

La première écriture $a = bq + r$ avec $r < b$ donne donc $q = \lfloor a/b \rfloor$ et
$r = a \bmod b$ ; la seconde donne de même $q' = \lfloor a/b \rfloor$ et
$r' = a \bmod b$. D'où $q = q'$ et $r = r'$.
\end{proof}

\source{https://github.com/Commutator-IO/learn-lean/blob/main/cours/01-college/07-enonces-conserves/EntiersDivisibilite.lean\#L169}{EntiersDivisibilite.lean\#L169}

\section{Nombre premier ; tout entier > 1 admet un diviseur premier}

\begin{theoreme}
Nombre premier : ses seuls diviseurs sont 1 et lui-même.
\end{theoreme}

\begin{proof}
C'est la seconde composante de la définition de \texttt{Premier}.
\end{proof}

\source{https://github.com/Commutator-IO/learn-lean/blob/main/cours/01-college/07-enonces-conserves/EntiersDivisibilite.lean\#L181}{EntiersDivisibilite.lean\#L181}

\begin{theoreme}
Tout entier > 1 admet un diviseur premier.
\end{theoreme}

\begin{proof}
Par récurrence forte sur $n$. Distinguons deux cas.

Si $n$ est premier, il est son propre diviseur premier.

Sinon, la négation de la définition fournit, par l'absurde, un diviseur $d$ de $n$ avec
$d \ne 1$ et $d \ne n$. Ce $d$ n'est pas nul : $0 \mid n$ entraînerait $n = 0$, contraire
à $n \ge 2$. Donc $2 \le d$. De $d \mid n$ et $n > 0$ on tire $d \le n$, et $d \ne n$
donne $d < n$. L'hypothèse de récurrence appliquée à $d$ fournit un premier $p$ divisant
$d$, et $p \mid n$ par transitivité.
\end{proof}

\source{https://github.com/Commutator-IO/learn-lean/blob/main/cours/01-college/07-enonces-conserves/EntiersDivisibilite.lean\#L185}{EntiersDivisibilite.lean\#L185}

\section{Crible d'Ératosthène : lister les nombres premiers inférieurs à 100}

\begin{definition}
Test calculable de primalité : $\textsf{estPremier}(n)$ vaut vrai lorsque $2 \le n$ et
que, pour tout $d < n$, on a $d < 2$ ou $n \bmod d \ne 0$. C'est le crible, écrit comme un
test.
\end{definition}

\source{https://github.com/Commutator-IO/learn-lean/blob/main/cours/01-college/07-enonces-conserves/EntiersDivisibilite.lean\#L211}{EntiersDivisibilite.lean\#L211}

\begin{theoreme}
Le test et la définition coïncident.
\end{theoreme}

\begin{proof}
($\Rightarrow$) Supposons le test satisfait et soit $d \mid n$. Comme $n \ge 2 > 0$, on a
$d \le n$. Si $d = n$, la conclusion est acquise. Si $d < n$, le test donne $d < 2$ ou
$n \bmod d \ne 0$ ; or $d \mid n$ force $n \bmod d = 0$, donc $d < 2$, c'est-à-dire
$d = 0$ ou $d = 1$. Le cas $d = 0$ est exclu, car $0 \mid n$ donnerait $n = 0$. Reste
$d = 1$.

($\Leftarrow$) Supposons $n$ premier et soit $d < n$. Si $d < 2$, la disjonction du test
est satisfaite. Sinon, supposons $n \bmod d = 0$ : alors $d \mid n$, donc $d = 1$ ou
$d = n$ par primalité, ce qui contredit $2 \le d < n$.
\end{proof}

\source{https://github.com/Commutator-IO/learn-lean/blob/main/cours/01-college/07-enonces-conserves/EntiersDivisibilite.lean\#L215}{EntiersDivisibilite.lean\#L215}

\begin{theoreme}
Crible d'Ératosthène : les nombres premiers inférieurs à 100.
\end{theoreme}

\begin{proof}
Par calcul : le noyau évalue la liste $0, 1, \dots, 99$, la filtre par le test et compare
le résultat à la liste annoncée. Le crible est donc exécuté, et pas seulement décrit —
c'est le théorème précédent qui garantit que ce que le test retient est bien la
primalité.
\end{proof}

\source{https://github.com/Commutator-IO/learn-lean/blob/main/cours/01-college/07-enonces-conserves/EntiersDivisibilite.lean\#L247}{EntiersDivisibilite.lean\#L247}

\section{Décomposition en produit de facteurs premiers}

\begin{definition}
Le \emph{produit} d'une liste vaut $1$ pour la liste vide, et $p \cdot \Pi(l)$ pour la
liste $p :: l$.
\end{definition}

\source{https://github.com/Commutator-IO/learn-lean/blob/main/cours/01-college/07-enonces-conserves/EntiersDivisibilite.lean\#L256}{EntiersDivisibilite.lean\#L256}

\begin{theoreme}
Décomposition en produit de facteurs premiers.
\end{theoreme}

\begin{proof}
Par récurrence forte sur $n$. Si $n = 1$, la liste vide convient, son produit valant $1$.

Sinon $n \ge 2$, et le théorème précédent fournit un premier $p$ divisant $n$, disons
$n = pm$. Alors $m \ne 0$, sinon $n$ serait nul ; et $m < n$, car
$m = 1 \cdot m < p \cdot m = n$ puisque $p \ge 2$ et $m > 0$. L'hypothèse de récurrence
donne une liste $l$ de premiers de produit $m$ ; la liste $p :: l$ convient alors pour
$n$ : tous ses éléments sont premiers, et son produit vaut $p \cdot m = n$.
\end{proof}

\source{https://github.com/Commutator-IO/learn-lean/blob/main/cours/01-college/07-enonces-conserves/EntiersDivisibilite.lean\#L261}{EntiersDivisibilite.lean\#L261}

\begin{lemme}
Lemme d'Euclide, sur lequel repose l'unicité.
\end{lemme}

\begin{proof}
Distinguons selon que $p$ divise $a$ ou non. S'il le divise, c'est fini.

Sinon, le pgcd $\gcd(p, a)$ divise $p$, donc vaut $1$ ou $p$ par primalité de $p$. S'il
valait $p$, alors $p$ diviserait $a$, ce qui est exclu. Donc $\gcd(p,a) = 1$ : $p$ et $a$
sont premiers entre eux, et de $p \mid ab$ on conclut $p \mid b$
(\texttt{Nat.Coprime.dvd\_of\_dvd\_mul\_left}).
\end{proof}

\source{https://github.com/Commutator-IO/learn-lean/blob/main/cours/01-college/07-enonces-conserves/EntiersDivisibilite.lean\#L290}{EntiersDivisibilite.lean\#L290}

\noindent
Statut ◐ pour cet énoncé : l'existence de la décomposition est démontrée ci-dessus, et
le lemme d'Euclide en donne la clé, mais l'unicité elle-même ne l'est pas encore. Elle
demande de comparer deux listes de facteurs à permutation près (\texttt{List.Perm}), ce qui est
une récurrence d'une tout autre longueur que l'énoncé de 3e ne le laisse croire \textemdash{} au
collège, l'unicité est admise sans même être formulée.

\section{PGCD, algorithme d'Euclide}

\begin{theoreme}
PGCD, algorithme d'Euclide : \texttt{pgcd(a, b) = pgcd(b, a mod b)}.
\end{theoreme}

\begin{proof}
C'est l'équation de récursion de \texttt{Nat.gcd} : la définition de la fonction
\emph{est} l'algorithme d'Euclide, et sa terminaison tient à la décroissance stricte de
$b \bmod a$.
\end{proof}

\source{https://github.com/Commutator-IO/learn-lean/blob/main/cours/01-college/07-enonces-conserves/EntiersDivisibilite.lean\#L309}{EntiersDivisibilite.lean\#L309}

\begin{theoreme}
Le PGCD est un diviseur commun.
\end{theoreme}

\begin{proof}
Conjonction de \texttt{Nat.gcd\_dvd\_left} et \texttt{Nat.gcd\_dvd\_right}.
\end{proof}

\source{https://github.com/Commutator-IO/learn-lean/blob/main/cours/01-college/07-enonces-conserves/EntiersDivisibilite.lean\#L313}{EntiersDivisibilite.lean\#L313}

\begin{theoreme}
Tout diviseur commun divise le PGCD.
\end{theoreme}

\begin{proof}
\texttt{Nat.dvd\_gcd}. Noter que la maximalité s'exprime ici par la divisibilité et non
par l'ordre : tout diviseur commun \emph{divise} le pgcd, ce qui est plus fort que de lui
être inférieur.
\end{proof}

\source{https://github.com/Commutator-IO/learn-lean/blob/main/cours/01-college/07-enonces-conserves/EntiersDivisibilite.lean\#L317}{EntiersDivisibilite.lean\#L317}

\section{Deux nombres sont premiers entre eux \(\iff\) \texttt{pgcd = 1}}

\begin{theoreme}
Nombres premiers entre eux \(\iff\) \texttt{pgcd = 1}.
\end{theoreme}

\begin{proof}
Vrai par définition : sur $\mathbb{N}$, « premiers entre eux » \emph{est} l'égalité du
pgcd à $1$.
\end{proof}

\source{https://github.com/Commutator-IO/learn-lean/blob/main/cours/01-college/07-enonces-conserves/EntiersDivisibilite.lean\#L324}{EntiersDivisibilite.lean\#L324}

\begin{theoreme}
Nombres premiers entre eux \(\iff\) seuls diviseurs communs égaux à 1.
\end{theoreme}

\begin{proof}
($\Rightarrow$) Un diviseur commun $d$ de $a$ et $b$ divise leur pgcd, qui vaut $1$ ;
donc $d = 1$.

($\Leftarrow$) Il suffit d'appliquer l'hypothèse au pgcd lui-même, qui est un diviseur
commun de $a$ et $b$.
\end{proof}

\source{https://github.com/Commutator-IO/learn-lean/blob/main/cours/01-college/07-enonces-conserves/EntiersDivisibilite.lean\#L328}{EntiersDivisibilite.lean\#L328}

\section{Toute fraction admet une écriture irréductible}

\begin{theoreme}
Toute fraction admet une écriture irréductible.
\end{theoreme}

\begin{proof}
Posons $g = \gcd(n,d)$, non nul puisque $d$ l'est, puis $n' = n/g$ et $d' = d/g$.

Alors $d' > 0$, car $g \le d$ ; et $n'$ et $d'$ sont premiers entre eux
(\texttt{Nat.coprime\_div\_gcd\_div\_gcd}). Enfin, de $g \cdot (n/g) = n$ et
$g \cdot (d/g) = d$ :
\[
n \cdot \frac{d}{g}
 = \Bigl(g \cdot \frac{n}{g}\Bigr) \cdot \frac{d}{g}
 = \frac{n}{g} \cdot \Bigl(g \cdot \frac{d}{g}\Bigr)
 = \frac{n}{g} \cdot d ,
\]
la deuxième égalité n'étant que la commutativité et l'associativité du produit. L'égalité
des fractions est écrite en produits croisés pour rester dans $\mathbb{N}$, où $n/d$ n'a
pas de sens.
\end{proof}

\source{https://github.com/Commutator-IO/learn-lean/blob/main/cours/01-college/07-enonces-conserves/EntiersDivisibilite.lean\#L341}{EntiersDivisibilite.lean\#L341}

\end{document}
